\scanpage{47}
\mspage{-3.4-}
% archive sheet number: 46
space. Then there is a $\mathrm{PGL}(n)$-linear isomorphism:
\begingroup\makeatletter\tagsleft@true\makeatother
\[
\handwritten{\unsure{\mathcal{U}|\mathscr{U}}}_R\handwritten{\cong}\mathrm{PGL}(n)\mathbin{\mathrm{x}}
\handwritten{\mathbb{A}}^{n\unsure{\nu|v}-n}\tag*{(*)}
\]
\endgroup
Hence $\handwritten{\mathcal{U}}_R$ is a globally trivial principal fibre bundle with respect
to the action of $\mathrm{PGL}(n)$, with base space $\handwritten{\mathbb{A}}^{n\nu-n}$.

\uline{Proof}: We shall construct morphisms from the left hand side of
$(*)$ to the right, and vice versa, leaving it to the reader to check
that they are mutually inverse. The morphisms will be obviously $\mathrm{PGL}(n)$-linear. First of all, we identify $\handwritten{\mathbb{A}}^{n\nu-n}$ with the scheme:
\[
\handwritten{\mathbb{A}}_R=\operatorname{Spec}
\frac{\handwritten{\mathbb{Z}}\bigl[x_0^{(n+1)},\ldots,x_n^{(n+1)};\ x_0^{(n+2)},\ldots,x_n^{(n+2)};\ldots;\ x_0^{(n+\nu)},\ldots,x_n^{(n+\nu)}\bigr]}
{\text{Ideal generated by }\handwritten{\{}x_i^{(n+k)}-1\mathrel{\big|}i\in S_k\handwritten{\}}}
\]
since both are affine spaces of dimension $n\nu-n$. Then a morphism
\[
\unsure{S|s}_R:\handwritten{\mathbb{A}}_R\handwritten{\longrightarrow}\handwritten{\mathcal{U}}_R
\]
is defined by the collection of composite morphisms
\[
S_i:\handwritten{\mathbb{A}}_R\handwritten{\xrightarrow{S}}\handwritten{\mathcal{U}}_R\handwritten{\subset}
(\handwritten{\mathbb{P}}_n)^{n+\nu+1}\handwritten{\xrightarrow{p_i}}\handwritten{\mathbb{P}}_n
\]
and $S_i$ is defined by the conditions (cf. EGA 2, \handwritten{4.2.3} ):
\[
\begin{aligned}
 S_i^*(\underline{o}\,(1))&\handwritten{\cong}\underline{o}_{\handwritten{\mathbb{A}}}\\
 S_i^*(X_j)&=\handwritten{\delta}_{ij}\quad\text{if }\ 0\leq i\leq n\\
 S_i^*(X_j)&=x_j^{(i)}\quad\text{if }\ n+1\leq i\leq m
\end{aligned}
\]

